% !TEX root =./ECE448_Practice_main.tex
%
% L8 practice set (FSK). Short by design: the L8 Lab packet carries the
% hands-on work. Blank copy: \begin{solution}[h] reserves h of work space.
% SOLUTIONS copy (\iskey=1) prints the red worked solutions.

\fullwidth{\textbf{Learning Objective 4.1: } I can describe FSK modulation and implement
an FSK transmitter and receiver in GNU Radio.
\\[4pt]\normalfont\small Show your work. Write a \textbf{Documentation} line at the
end (\textbf{None} if you did not collaborate).}

\begin{questions}

%%%%%%%%%%%%%%%%%%%%%%%%%%%%%%%%%%%%%%%%%
\question \textbf{Transmitter settings.} A flowgraph samples at $f_s = 48$ kHz and
sends FSK at 1200 baud with a deviation of $\Delta f = 600$ Hz, one bit per symbol.
\begin{parts}

	\begin{minipage}[t]{0.58\linewidth}
		\part Find the samples per symbol to enter in the Repeat block.
		\begin{solution}[0.7in]
		\eq{ \text{sps} = \dfrac{f_s}{R_s} = \dfrac{48000}{1200} = 40 }
		\end{solution}
	\end{minipage}\hfill%
	\begin{minipage}[t]{0.37\linewidth}
		\raggedleft \vspace{12pt}
		\ansbox[1.3in]{40}
	\end{minipage}

	\begin{minipage}[t]{0.58\linewidth}
		\part Find the sensitivity to enter in the Frequency Mod block.
		\begin{solution}[0.8in]
		\eq{ k = \dfrac{2\pi \Delta f}{f_s} = \dfrac{2\pi \times 600}{48000} = 0.0785\ \text{rad/sample} }
		\end{solution}
	\end{minipage}\hfill%
	\begin{minipage}[t]{0.37\linewidth}
		\raggedleft \vspace{12pt}
		\ansbox[1.3in]{0.0785}
	\end{minipage}

	\begin{minipage}[t]{0.58\linewidth}
		\part Find the modulation index and estimate the bandwidth with Carson's rule.
		\begin{solution}[1.0in]
		\eq{ h &= \dfrac{2\Delta f}{R_s} = \dfrac{1200}{1200} = 1 \\
			B &\approx 2(\Delta f + R_s) = 2(600 + 1200) = 3.6\ \text{kHz} }
		\end{solution}
	\end{minipage}\hfill%
	\begin{minipage}[t]{0.37\linewidth}
		\raggedleft \vspace{12pt}
		$h$ = \ansbox[1.0in]{1}

		\vspace{4pt}
		$B$ = \ansbox[1.0in]{3.6 kHz}
	\end{minipage}

	\part A classmate builds the same transmitter from two Signal Sources at $+600$
	and $-600$ Hz and a selector that switches between them at each bit. Explain why
	Frequency Mod gives a more compact spectrum.
		\begin{solution}[1.0in]
		The two sources run independently, so at each switch the phase jumps from
		wherever one source is to wherever the other is. Those jumps are abrupt
		changes in the signal, and abrupt changes spread power across a wide band.
		Frequency Mod adds a phase step every sample, so the phase is continuous
		across each bit boundary and the spectrum stays compact.
		\end{solution}

\end{parts}

\newpage
%%%%%%%%%%%%%%%%%%%%%%%%%%%%%%%%%%%%%%%%%
\question \textbf{Receiver settings.} The receiver for the link of Question 1 follows
the L8 flowgraph: Low Pass Filter, Quadrature Demod, Moving Average, Symbol Sync,
and Binary Slicer.
\begin{parts}

	\begin{minipage}[t]{0.58\linewidth}
		\part Find the Quadrature Demod gain that puts the two tones at $+1$ and $-1$.
		\begin{solution}[0.8in]
		\eq{ g = \dfrac{f_s}{2\pi\Delta f} = \dfrac{48000}{2\pi \times 600} = 12.7 }
		\end{solution}
	\end{minipage}\hfill%
	\begin{minipage}[t]{0.37\linewidth}
		\raggedleft \vspace{12pt}
		\ansbox[1.3in]{12.7}
	\end{minipage}

	\begin{minipage}[t]{0.58\linewidth}
		\part Find the low-pass filter cutoff that passes the Carson bandwidth.
		\begin{solution}[0.7in]
		The filter passes $-f_c$ to $+f_c$, so the cutoff is half of $B$:
		\eq{ f_c = \Delta f + R_s = 600 + 1200 = 1800\ \text{Hz} }
		\end{solution}
	\end{minipage}\hfill%
	\begin{minipage}[t]{0.37\linewidth}
		\raggedleft \vspace{12pt}
		\ansbox[1.3in]{1800 Hz}
	\end{minipage}

	\part A classmate leaves the Quadrature Demod gain at 1. Find the levels the two
	tones produce, and explain whether the slicer still decides correctly.
		\begin{solution}[1.1in]
		With $g = 1$ the output is the phase step itself,
		$\pm 2\pi\Delta f/f_s = \pm 0.0785$. The slicer compares with zero, and both
		levels are still on the correct side of zero, so the bits come out correctly.
		Only the scale is wrong: the eye diagram and any plot set for $\pm 1$ show a
		nearly flat line.
		\end{solution}

	\begin{minipage}[t]{0.58\linewidth}
		\part The classmate also sets the filter cutoff to 8 kHz. By how many dB does
		that lower the SNR at the demodulator?
		\begin{solution}[0.9in]
		The filter now passes 16 kHz of noise instead of 3.6 kHz, and the signal is
		unchanged:
		\eq{ 10\mylog{\dfrac{16}{3.6}} = 6.5\db }
		\end{solution}
	\end{minipage}\hfill%
	\begin{minipage}[t]{0.37\linewidth}
		\raggedleft \vspace{12pt}
		\ansbox[1.3in]{$6.5\db$}
	\end{minipage}

\end{parts}

\newpage
%%%%%%%%%%%%%%%%%%%%%%%%%%%%%%%%%%%%%%%%%
\question \textbf{Noise.} The L8 link samples at 32 kHz, sends 1000 baud with
$\Delta f = 1$ kHz, and filters to the 4 kHz Carson bandwidth. The signal has power
1, and the Channel Model's noise of voltage $\sigma$ has power $\sigma^2$ spread over
the full 32 kHz.
\begin{parts}

	\begin{minipage}[t]{0.58\linewidth}
		\part Find the SNR at the demodulator for $\sigma = 0.8$.
		\begin{solution}[0.9in]
		\eq{ \text{SNR} = \dfrac{1}{\sigma^2 \times 4/32} = \dfrac{8}{0.64} = 12.5 = 11.0\db }
		\end{solution}
	\end{minipage}\hfill%
	\begin{minipage}[t]{0.37\linewidth}
		\raggedleft \vspace{12pt}
		\ansbox[1.3in]{$11.0\db$}
	\end{minipage}

	\begin{minipage}[t]{0.58\linewidth}
		\part The text starts to show errors below about 9 dB. Find the largest noise
		voltage that keeps the SNR at 9 dB.
		\begin{solution}[0.9in]
		\eq{ \sigma^2 = \dfrac{8}{10^{0.9}} = 1.01, \qquad \sigma = 1.0 }
		\end{solution}
	\end{minipage}\hfill%
	\begin{minipage}[t]{0.37\linewidth}
		\raggedleft \vspace{12pt}
		\ansbox[1.3in]{$\sigma = 1.0$}
	\end{minipage}

	\begin{minipage}[t]{0.58\linewidth}
		\part The link doubles its symbol rate to 2000 baud and keeps $\Delta f = 1$ kHz.
		Find the new Carson bandwidth and the largest noise voltage for 9 dB.
		\begin{solution}[1.1in]
		\eq{ B &= 2(1000 + 2000) = 6\ \text{kHz} \\
			\sigma^2 &= \dfrac{32/6}{10^{0.9}} = 0.67, \qquad \sigma = 0.82 }
		\end{solution}
	\end{minipage}\hfill%
	\begin{minipage}[t]{0.37\linewidth}
		\raggedleft \vspace{12pt}
		$B$ = \ansbox[1.0in]{6 kHz}

		\vspace{4pt}
		$\sigma$ = \ansbox[1.0in]{0.82}
	\end{minipage}

	\part Lesson 2 designs a link with margin above the SNR it needs. Explain, from
	what the L8 link does between 9 dB and 3 dB, why margin matters for FSK.
		\begin{solution}[1.0in]
		The link goes from almost no errors at 9 dB to almost no packets at 3 dB, a
		range of about 6 dB. A link designed exactly at 9 dB therefore fails completely
		after a fade of a few dB. Margin keeps the SNR above the threshold when the
		path loss changes.
		\end{solution}

\end{parts}

\newpage
%%%%%%%%%%%%%%%%%%%%%%%%%%%%%%%%%%%%%%%%%
\question \textbf{An FSK remote on a waterfall.} An RTL-SDR waterfall shows a burst
with two tracks, at 433.900 and 433.940 MHz. In the alternating preamble, the signal
stays on each track for 0.417 ms.
\begin{parts}

	\begin{minipage}[t]{0.58\linewidth}
		\part Find the center frequency, the deviation, and the symbol rate.
		\begin{solution}[1.1in]
		\eq{ f_c &= \dfrac{433.900 + 433.940}{2} = 433.920\ \text{MHz} \\
			\Delta f &= \dfrac{40\ \text{kHz}}{2} = 20\ \text{kHz} \\
			R_s &= \dfrac{1}{0.417\ \text{ms}} = 2400\ \text{baud} }
		\end{solution}
	\end{minipage}\hfill%
	\begin{minipage}[t]{0.37\linewidth}
		\raggedleft \vspace{12pt}
		$f_c$ = \ansbox[1.2in]{433.920 MHz}

		\vspace{4pt}
		$\Delta f$ = \ansbox[1.2in]{20 kHz}

		\vspace{4pt}
		$R_s$ = \ansbox[1.2in]{2400 baud}
	\end{minipage}

	\begin{minipage}[t]{0.58\linewidth}
		\part The remote's crystal is accurate to 10 ppm. Find the largest frequency
		offset at 433.92 MHz.
		\begin{solution}[0.7in]
		\eq{ 10\times10^{-6} \times 433.92\times10^{6} = 4.3\ \text{kHz} }
		\end{solution}
	\end{minipage}\hfill%
	\begin{minipage}[t]{0.37\linewidth}
		\raggedleft \vspace{12pt}
		\ansbox[1.3in]{4.3 kHz}
	\end{minipage}

	\part An FSK receiver tolerates an offset of about half its deviation. Explain
	whether a 4.3 kHz offset breaks the remote's link, and whether it would break the
	L8 link if that link ran at 433.92 MHz.
		\begin{solution}[1.1in]
		The remote tolerates about 10 kHz, so 4.3 kHz is well inside its limit; the
		large deviation is what makes a cheap crystal acceptable. The L8 link tolerates
		about 600 Hz, so a 4.3 kHz offset pushes both of its tones far off the
		receiver's filter and the link fails.
		\end{solution}

	\part With a 600 Hz offset the L8 demodulator outputs levels near $+1.6$ and
	$-0.4$. Name one change to the receiver that would put them back at $\pm 1$.
		\begin{solution}[0.9in]
		Subtract the average of the demodulator output, for example with a DC Blocker
		after the Quadrature Demod, because the offset appears as a constant
		$f_\text{off}/\Delta f$ added to both levels. Measuring the offset and shifting
		the signal back in frequency before the filter also works.
		\end{solution}

\end{parts}

\vspace{1.5em}
\noindent\textbf{Documentation:} \rule[-2pt]{0.6\linewidth}{0.4pt}

\end{questions}
